\par\noindent\textbf{(1)} 单粒子能量 \(\varepsilon=p^2/(2m)\)，激发态平均数 \(N_{\rm ex}=\frac{V}{(2\pi\hbar)^3}\int d^3p[\exp(\beta(\varepsilon-\mu))-1]^{-1}\)。令热德布罗意波长 \(\lambda_T=\sqrt{2\pi\hbar^2/(mk_BT)}\)、逸度 \(z=e^{\beta\mu}\le1\)，则 \(N_{\rm ex}=V\lambda_T^{-3}g_{3/2}(z)\)。
\par\noindent\textbf{(2)} 当 \(\mu\uparrow0\)，激发数上限 \(N_{\rm ex,max}=V\lambda_T^{-3}\zeta(3/2)\)。定密度 \(n=N/V\) 下 \(T_c=(2\pi\hbar^2/(mk_B))[n/\zeta(3/2)]^{2/3}\)。若 \(T<T_c\)，剩余占基态 \(N_0/N=1-(T/T_c)^{3/2}\)。
\par\noindent\textbf{(3)} 在 \(\mu=0\) 的凝聚区，激发态内能 \(U=\tfrac32 Vk_BT\lambda_T^{-3}\zeta(5/2)\propto T^{5/2}\)，所以 \(C_V=\partial_TU=\tfrac{15}{4}Vk_B\lambda_T^{-3}\zeta(5/2)\)。基态能零、宏观占有但不贡献理想气体热能。
\par\noindent\textbf{(4)} 对宏观占据的单基态，巨正则分布为几何分布，\(\operatorname{Var}(N_0)=\bar N_0(1+\bar N_0)\)，可达到 \(O(N^2)\) 的异常涨落。固定 \(N\) 的正则系综不存在此独立几何涨落，因此用巨正则凝聚态结果外推真实封闭气体要注意系综效应。
